\documentclass[PhD,single]{uclathes}

\usepackage{graphicx}
\usepackage{amsmath}
\usepackage{amssymb}
\usepackage[ruled,vlined,linesnumbered]{algorithm2e}

\title{Graph Neural Networks on Large Graphs via Graph Databases}
\author{Victor Pekkari, Harry Wik}
\university{Lund University, Faculty of Engineering\\Neo4j}
\department{Computer Science}
\thesis{Master of Science}
\degreemonth{March}
\degreeyear{2026}
\nocopyright

\acknowledgments{%
We would like to thank the Graph Data Science team at Neo4j for hosting this
thesis work and for providing access to infrastructure, expertise, and an
invaluable working environment. We are also grateful to our supervisor and
examiner at LTH for their guidance throughout the project.
}

\abstract{% 
Graph-structured data appears in domains such as fraud detection, recommendation,
and molecular analysis, which has made graph neural networks (GNNs) a central
approach for learning on graphs. However, message passing creates dependencies
between nodes that make large-scale training and inference difficult, especially
when the graph does not fit comfortably in memory. Existing large-scale GNN
systems commonly rely on distributed training and graph partitioning, but these
approaches can suffer from communication overhead and expensive sampling.

This thesis investigates whether a graph database can serve as an efficient and
scalable backend for GNN training and inference. More specifically, it studies a
Neo4j-backed pipeline integrated with PyTorch Geometric and evaluates the
bottlenecks of database-backed training, the efficiency of different sampling
strategies, the best approach for inference on new nodes, and how the usefulness
of the approach changes with graph size. The work follows an experimental,
systems-oriented methodology where components of a functionally equivalent
PyTorch Geometric pipeline are replaced, benchmarked, and analyzed.

The thesis contributes a structured system design for graph-database-backed GNN
training, a methodology for evaluating both correctness and throughput, and an
analysis framework for understanding when graph databases are competitive with
more conventional large-scale GNN systems.
}

\begin{document}

\makeintropages

\chapter{Introduction}

Graph structures are a natural way to model entities and their relationships in
several important domains. Such models are applied in diverse domains, including
predicting chemical bonds in medical science and detecting fraud in finance. The
prevalence of graph-structured data across many domains creates a need for
efficient and accurate methods to model and analyze graphs.

Graph neural networks (GNNs) have emerged as a leading approach for
representation learning on graphs, combining the structural intuition behind
classical graph algorithms with the expressiveness and inductive capabilities of
neural networks.

However, GNNs introduce a challenge absent in traditional neural networks:
message passing creates growing dependencies between nodes during training and
inference.

Existing approaches for scaling GNNs to large graphs rely on partitioning the
graph, sampling subgraphs, and distributing computation across machines. These
approaches often suffer from high communication overhead between machines, and
from inefficient sampling procedures that can take longer than the actual model
training.

In practice, many large graphs are already stored in graph databases before they
are loaded into workers for GNN training or inference. The characteristics of
fast graph traversal and practically unlimited memory in graph databases make
them a promising component in a GNN training and inference pipeline for large
graphs. While it has been shown that it is possible to use a graph database as a
backend for GNN training, prior work has mainly demonstrated feasibility rather
than performance competitiveness with state-of-the-art distributed systems. It
is therefore still unclear (i) what the bottlenecks in graph-database-backed
training and inference are, (ii) what an efficient graph-database-based pipeline
would look like, (iii) whether and when graph-database-backed training
pipelines might be a better choice than distributed systems, and (iv) how
inference should be done with the graph stored in a graph database.

This thesis investigates whether a graph database can serve as an efficient and
scalable backend for GNN training and inference on large graphs. To answer this,
we implement and evaluate a Neo4j-backed pipeline integrated with PyTorch
Geometric, analyze its main bottlenecks, compare sampling strategies, and study
how the usefulness of this approach depends on graph size and inference setting.

\section{Research questions}
This thesis addresses the following research questions:

\begin{enumerate}
    \item \textbf{RQ1 --- Bottlenecks:}\label{question:bottleneck} What are the primary bottlenecks in GNN training when using a graph database for graph storage, and to what extent can they be mitigated?
    \item \textbf{RQ2 --- Sampling:}\label{question:sampling} Which subgraph sampling strategies are most efficient when the graph is stored in a graph database?
    \item \textbf{RQ3 --- Inference:}\label{question:inference} What is the optimal method for doing inference on new nodes when the graph is stored in a graph database?
    \item \textbf{RQ4 --- Graph size:}\label{question:graphsize} For which graph sizes does a graph database-based storage approach provide the greatest benefits?
\end{enumerate}

\section{Scope and limitations}
The experiments in this thesis are conducted on four datasets of varying sizes
using a two-layer Graph Convolutional Network (GCN) for node classification.
The main emphasis is on the training and inference pipeline rather than on model
architecture. The study is therefore limited to a specific GNN setting, but
targets system-level questions that are relevant to large-scale graph learning
more broadly.

\section{Thesis outline}
Chapter 2 provides background on important topics needed for understanding the
rest of the thesis, such as GNNs, neural networks, graph representation
learning, Neo4j, and PyTorch Geometric. Chapter 3 reviews related work on GNNs
applied to large graphs. Chapter 4 describes the methodology used during the
experiments. Chapter 5 outlines the implemented GNN pipelines. Chapter 6
describes the results, Chapter 7 discusses them, and Chapter 8 concludes the
thesis.

\chapter{Background}

Graphs in mathematics are used to model pairwise relationships between objects.
Objects are modeled as nodes $\mathcal{V}$, and their relations are called
edges $\mathcal{E}$. Many domains benefit from modeling data as objects (nodes)
and relationships (edges). Consider, for example, social networks,
protein--drug interactions, or railway networks. Studying graph structure and
node features allows us to predict properties of the graph.

When making predictions about graphs it is often useful to start by creating an
embedding of the graph. The optimal embedding depends on the downstream task.
There are several ways to embed a graph, transductively with traditional graph
representation learning (GRL) approaches, and inductively with GNNs.

The purpose of this chapter is to give a fundamental understanding of
traditional GRL algorithms, neural networks, and how GNNs combine ideas from
traditional graph representation learning and neural networks. From traditional
GRL, GNNs inherit the idea that node embeddings should reflect graph structure,
while from neural networks they adopt a learnable, parameterized function.
Challenges of working with GNNs on large datasets that do not exist in
traditional neural networks are also discussed.

Lastly, a basic introduction to \texttt{PyTorch Geometric}, the machine
learning library used during experiments, and \texttt{Neo4j}, a graph database
that provides useful graph operations for training GNNs, is provided.

\section{Graphs and graph representation learning}\label{sec:traditional_grl}
In graph representation learning the aim is to find a representation of the
graph that encodes its structure. Embeddings like this are often used in
downstream tasks such as node classification, link prediction, or clustering.
Finding a good embedding depends in part on what the downstream task is.

\subsection{Transductive nature of traditional GRL}
A common feature throughout traditional graph representation learning
algorithms, whether based on node statistics, matrix factorization, or random
walks, is that no weights are shared between nodes and no node-embedding
generating function is learned. Instead, only node embeddings for the current
nodes in the graph are learned. This means that if new nodes are introduced to
the graph, there is no direct way to compute their embeddings, and the
embeddings for the old nodes are no longer accurate. The latter also happens if
nodes disappear from the graph.

Consider a simplification of the optimization problem for random-walk-based
algorithms:

\[
\min_{\{z_x \,\,: \,\,x \in \mathcal{V} \}}
\sum_{(u,v)\in\mathcal{D}}
-\log\left(\sigma(\mathbf{z}_u^\top \mathbf{z}_v)\right)
\]

where $\mathcal{D}$ is a set of node pairs encountered on the same random walk,
and $\mathcal{V}$ is the set of nodes in the graph.

The optimization problem above is an example of how it is the actual node
embeddings that are optimized, rather than some embedding-generating function.
This is what makes traditional representation learning algorithms transductive
in nature.

Consider a social network where people are nodes, and their relations are
edges. If a transductive graph representation is used, the learned feature
vectors are tied to specific people in the graph. If more people are introduced
to the network, the feature vector for every person has to be recomputed.

\section{Neural networks}
Neural networks are parameterized functions that can learn a function mapping
inputs to outputs. Given an input $x$, a neural network produces an output
$\hat{y}$ according to a function $f_{\theta}(x)$, where $\theta$ are the
parameters, also known as weights, of the model.

Optimizing the weights $\theta$ of the model includes minimizing a loss
function $\mathcal{L}$ which measures how far the model prediction is from the
desired target values. This optimization is usually done using some
gradient-based method.

A key property of neural networks is that they are inductive: instead of
learning representations tied to specific data points, they learn a general
function that can be applied to unseen inputs. This enables generalization
beyond the training data.

Standard neural network training further relies on the assumption that data
points are independent. Under this assumption, the loss can be approximated
using mini-batches of data, allowing efficient and scalable training. However,
this assumption does not hold for graph-structured data, where data points
(nodes) are explicitly dependent on each other through the graph topology. This
limitation motivates the development of graph neural networks, which extend
neural networks to account for relational dependencies.

\section{Graph neural networks}
Graph neural networks (GNNs) are neural networks designed for graph-structured
data. Like traditional graph representation learning methods, they aim to
produce node representations that reflect graph structure. Unlike shallow graph
embedding methods, however, GNNs learn a shared parameterized function that
generates node representations from node features and graph connectivity. This
makes GNNs inductive: once trained, the same model can be applied to unseen
nodes, and in some cases to entirely new graphs drawn from a similar
distribution.

A central idea in GNNs is \emph{message passing}. Instead of computing a node
representation only from the node's own features, a GNN updates each node by
aggregating information from its neighbors. By stacking several such layers, the
model can incorporate information from increasingly distant parts of the graph.

\subsection{Message passing}
At message passing layer $k$, the representation of node $v$ is updated as
\[
\mathbf{x}_v^{(k)} =
\phi^{(k)}\!\left(
\mathbf{x}_v^{(k-1)},
\bigoplus_{u \in \mathcal{N}(v)}
\psi^{(k)}\!\left(\mathbf{x}_v^{(k-1)}, \mathbf{x}_u^{(k-1)}\right)
\right),
\]
where $\mathcal{N}(v)$ denotes the neighbors of node $v$, $\psi^{(k)}$ is a
message function, $\bigoplus$ is a permutation-invariant aggregation operator
such as sum or mean, and $\phi^{(k)}$ is an update function.

After $K$ layers, the representation of a node depends on information from its
$K$-hop neighborhood. This allows GNNs to exploit both node features and graph
structure when solving downstream tasks such as node classification.

\subsection{Implications for large-graph training}
The message-passing mechanism also creates the main computational challenge in
GNN training. In standard neural networks, data points in a mini-batch are
usually treated as independent, so the loss and gradient can be computed from
the batch alone. In a GNN, this is no longer true: the representation of a node
depends on its neighbors, whose representations in turn depend on their
neighbors. As a result, computing the output for a batch of seed nodes
generally requires access to a larger induced subgraph around those nodes.

This dependency grows with the number of GNN layers. For a $K$-layer model,
each seed node may depend on a $K$-hop neighborhood, which can become very
large in real-world graphs. Training on the full graph is therefore often
infeasible for memory and computation reasons, especially for large graphs. To
address this, practical GNN pipelines typically rely on subgraph or
neighborhood sampling, where only a bounded part of the neighborhood is
retrieved for each mini-batch.

\section{Graph neural networks for large graphs}
In typical neural networks, data points are either independent or have
dependencies with a clear and structured form. This structure is determined by
the architectural inductive bias, which makes it straightforward to derive
which data points influence each other. For example, in recurrent neural
networks, $t_n$ depends on $\{t_{n-1}, t_{n-2}, \dots\}$, while in
convolutional neural networks a pixel $p_{x,y}$ depends only on a local
neighborhood $\{p_{i,j} \mid |i-x| \le r,\ |j-y| \le r\}$ in the same image.
In both cases, the set of dependent data points is explicitly defined and
bounded.

In contrast, graph neural networks lack such explicit structural assumptions.
The dependency structure is determined by the graph itself, making it harder to
determine which data points influence a given node. Moreover, these
dependencies are not only less structured, but also significantly larger. The
representation of a node depends on its $k$-hop neighborhood, which can grow
exponentially with $k$. As a result, each data point in a GNN may depend on
orders of magnitude more data points than in models such as convolutional
neural networks, where dependencies are limited to small local regions.

\subsection{Training under memory constraints}
In standard neural networks, the independence, or structured bounded
dependence, between data points allows training on datasets larger than the
available memory. Data can be processed in virtual mini-batches
$\{d^{(1)}, \dots, d^{(n)}\}$ such that
\[
\sum_i d_x^{(i)} < Y,
\]
where $Y$ is the available memory. Crucially, the gradient computed on a batch
depends only on the data points within that batch. This assumption breaks for
GNNs as the output for one node is dependent on an unbounded number of
neighbors. The gradient for the seed nodes in the batch cannot be computed
accurately without their neighbors. To solve this, the neighbors of the seed
nodes are fetched at the beginning of every batch to build a subgraph around the
seed nodes so that the gradient for that batch can be computed accurately.

\begin{algorithm}[H]
\label{algo:sampling}
\footnotesize
\caption{Mini-batch GNN Training (Single Epoch)}
\DontPrintSemicolon
\tcp{Mini-batch training with neighbor-sampled subgraphs}
\KwIn{$f_\theta$, $V$}
\KwOut{$\theta$}
\For{each batch $B \subseteq V$}{
    $\mathcal{G}_B \leftarrow \emptyset$\;
    \For{each seed node $v \in B$}{
        $(\mathbf{X}_v, \mathbf{y}_v, E_v) \leftarrow \textsc{Sample}(v)$\;
        $\mathcal{G}_B \leftarrow \mathcal{G}_B \cup (\mathbf{X}_v, \mathbf{y}_v, E_v)$\;
    }
    $\hat{\mathbf{y}} \leftarrow f_\theta(\mathbf{X}_B,\, E_B)$\;
    $\mathcal{L} \leftarrow \textsc{CrossEntropy}(\hat{\mathbf{y}},\, \mathbf{y}_B)$\;
    $\theta \leftarrow \theta - \eta \nabla_\theta \mathcal{L}$\;
}
\Return $\theta$\;
\end{algorithm}

\subsection{Subgraph sampling}
To address this, GNN training typically relies on subgraph sampling. Instead of
using full neighborhoods, a bounded subset of neighbors is sampled for each
node. Given a batch of seed nodes, a subgraph is constructed by sampling up to
a fixed number of neighbors at each hop. This limits the size of the
computation graph and enables mini-batch training.

A common approach is layer-wise neighbor sampling, where upper bounds are
defined as \texttt{[$k_0, k_1, \ldots, k_{K-1}$]}. At layer $i$, at most $k_i$
neighbors are sampled per node. If the number of neighbors exceeds this bound,
a subset is selected, typically uniformly at random.

While sampling makes training tractable, it introduces an approximation of the
true neighborhood aggregation. Consider the exact aggregation for a node $v$,
using the mean of node-neighbor feature vectors as the aggregation function:
\[
h_v = \sigma\!\left( \frac{1}{|\mathcal{N}(v)|} \sum_{u \in \mathcal{N}(v)} x_u \right),
\]
where $\sigma$ is a nonlinear activation function.

With sampling, we instead compute:
\[
\tilde{h}_v = \sigma\!\left( \frac{1}{k} \sum_{u \in \tilde{\mathcal{N}}(v)} x_u \right),
\]
where $\tilde{\mathcal{N}}(v)$ is a sampled subset of neighbors.

Although the sampled mean is an unbiased estimator of the true mean,
\[
\mathbb{E}\!\left[ \frac{1}{k} \sum_{u \in \tilde{\mathcal{N}}(v)} x_u \right]
=
\frac{1}{|\mathcal{N}(v)|} \sum_{u \in \mathcal{N}(v)} x_u,
\]
the presence of the nonlinearity introduces bias:
\[
\mathbb{E}[\tilde{h}_v]
=
\mathbb{E}\!\left[\sigma\!\left( \cdot \right)\right]
\neq
\sigma\!\left( \mathbb{E}[\cdot] \right)
=
h_v.
\]

This discrepancy arises because $\sigma$ is nonlinear, and therefore
expectation does not commute with the activation function. In practice, this
bias is often small and does not prevent convergence, especially for shallow
GNNs.

\section{PyTorch Geometric}\label{sec:pyg}
PyTorch Geometric is a graph machine learning library built on top of PyTorch.
It provides implementations of various message-passing layers, as well as
abstractions for storing the graph.

\begin{itemize}
    \item \texttt{FeatureStore:} defines how the node features are stored and retrieved.
    \item \texttt{GraphStore:} defines how the graph is traversed, and how edges are stored.
    \item \texttt{BaseSampler:} defines how neighborhood sampling is done, and is tightly coupled with the graph store.
    \item \texttt{NodeLoader:} takes implementations of \texttt{FeatureStore}, \texttt{GraphStore}, and \texttt{BaseSampler} and provides a data loader.
\end{itemize}

\section{Neo4j and graph databases}
Databases are systems used to efficiently store and fetch data. Instead of
scanning raw files, databases provide a way to find relevant data with low
latency. Traditional databases are relational, and the data is stored in the
form of tables of rows and columns. Relationships between objects are stored
through foreign keys, which leads to expensive join operations when retrieving
relational data. This model is not well suited for queries that involve
repeated traversal of relationships, which is something needed for fast
subgraph sampling in GNN training.

Graph databases such as Neo4j are, in contrast, optimized to have efficient
operations on graphs. This is done by explicitly modeling nodes and edges,
making relationships first-class citizens. This allows queries that follow
connections, such as multi-hop traversals, to be expressed and executed more
naturally and efficiently. In an ETL pipeline to a GNN model it is necessary
that the database supports efficient retrieval, and scalability in the form of
horizontal scaling via clustering as well as vertical scale-up of the machines
running the database instances.

\subsection{Alternatives to Neo4j}
Although Neo4j is a market leader, capturing the largest market share in the
graph database space, there are several alternatives. Amazon Neptune, Azure
Cosmos, and Google Spanner Graph are examples that mainly compete with Neo4j by
tight integration with the respective cloud platforms. Other examples include
Memgraph and TigerGraph.

\subsection{Querying Neo4j with Cypher and procedures}
Cypher is a declarative query language for graph databases where queries are
expressed as pattern matches over nodes and edges. The query planner translates
the Cypher code into execution plans optimized for traversal and lookup. In
this work, Cypher is used to express neighborhood expansion queries required
for subgraph sampling.

Procedures extend Cypher with imperative logic that cannot be expressed in
Cypher. Procedures are defined in Java within the database engine, allowing
custom traversal strategies, early termination, and more fine-grained control
over execution. This makes procedures a relevant alternative when trying to
optimize highly specialized Cypher queries. Procedures are explored alongside
Cypher when fetching neighbors of the seed nodes.

\subsection{Data retrieval from Neo4j}
Data is retrieved in batches determined by the \texttt{fetch\_size}. A smaller
\texttt{fetch\_size} results in more round trips between the client and the
database. For each batch, a number of rows corresponding to the
\texttt{fetch\_size} are first located in the database and then serialized
before being transmitted over the network. The serialization speed depends on
the underlying data format; for example, byte arrays are typically faster to
serialize than more complex structures. While one batch is being transmitted,
the next batch is prepared in the database, serialized, and subsequently sent
to the client.

\subsection{Cypher query execution pipeline}
The incoming Cypher string is parsed into an abstract syntax tree. The syntax
tree is then validated, for example by checking for labels, relationship types,
or property predicates that do not exist. The next step is normalization into a
canonical form.

The normalized syntax tree is used to generate one or more logical plans:
operator trees that define how to produce the query response. Much like the
expression
\begin{align*}
    R := c + a \cdot b
\end{align*}
is reduced according to operator precedence, a logical plan generated by a
Cypher query acts upon the data source with operators such as
\texttt{NodeByLabelScan} and \texttt{Filter}. One logical plan can often be
resolved according to many different physical plans, all yielding the same
output given the input. A simpler query is likely to only have one physical
plan whereas a more complex query can have many. In the case of multiple
physical plans, they are compared against each other so that a cost metric is
minimized. Examples of elements weighed into the cost optimization are metadata
such as how many nodes have a certain label and whether an index lookup should
be used.

After cost optimization one physical plan is promoted to be the execution plan.
The execution step compiles optimized bytecode that can be used and reused. To
answer the query in a timely manner, the execution step can have a parallel
runtime.

Some features are restricted to paying database clients, that is, clients that
run Neo4j Enterprise rather than Neo4j Community. Parallelism is one such
example along with clustering, extra constraints, online backups, and role
based access control.

Regardless of parallelism, execution first attempts to retrieve from the page
cache in RAM. The nodes and relations required to resolve the query are either
found in cache through pointer chasing, or a page fault occurs and data needs
to be loaded from non-volatile memory, for example an NVMe RAID disk. The exact
method of retrieval depends on the storage format. The current default is block
format.

After data has been localized, it needs to be sent back to the client. First,
the binary representations are cast to proper types. Next is the removal of
non-return variables. Before finally returning data to the client, the Java
classes must be serialized. Neo4j has a Bolt server that utilizes PackStream to
binary encode the data, after which the standard TCP/IP stack takes over and
answers the query.

\chapter{Related Work}

Existing approaches to scaling GNN training can be broadly categorized as
sampling-based methods and distributed systems. Sampling methods reduce
computational cost by approximating neighborhoods, while distributed systems
improve scalability through parallelism at the cost of communication overhead.

\section{Subgraph sampling}
Sampling methods are widely used when training GNNs because of the exploding
neighborhoods that come from stacked message-passing layers. Sampling-based
methods sample subgraphs from the larger original graph. This reduces both the
computational and memory cost.

\textbf{GraphSAGE} introduced layer-wise neighbor sampling, where a fixed
number of neighbors are sampled per node during training. \textbf{GraphSAINT}
instead samples subgraphs using random walks or node-based sampling strategies,
aiming to improve training efficiency and reduce variance.

While these methods make large-scale GNN training tractable, they introduce
approximation error due to incomplete neighborhood information, and their
performance depends on the efficiency of the sampling procedure.

\section{Distributed GNN systems}
Another line of work focuses on distributed GNN training, where the graph is
partitioned across multiple machines. These systems aim to scale training by
parallelizing both computation and data access.

[insert about papers here]

Despite their scalability, distributed approaches often suffer from
communication overhead between machines and require careful partitioning
strategies to balance load and minimize data transfer.

\section{Database-backed GNN pipelines}
A smaller body of work explores graph-database-backed GNN training. Here the
memory storage and fast graph traversal in graph databases is leveraged,
instead of storing the graph in memory across several machines.

Lopushanskyy and Shi demonstrate that it is possible to train GNNs directly on a
graph database by retrieving only the required subgraphs during training.

However, existing work in this area has primarily focused on demonstrating
feasibility, with limited evaluation of performance compared to state-of-the-
art distributed systems. In particular, it remains unclear what the main
bottlenecks are and under what conditions database-backed approaches are
competitive.

\section{Research gap}
The literature shows that both sampling-based and distributed approaches can
scale GNN training, but it leaves open several questions about graph-database-
backed pipelines. In particular, prior work does not sufficiently explain what
their dominant bottlenecks are, which sampling strategies are most effective in
this setting, how inference should be performed when the graph remains in the
database, or for which graph sizes such an approach is most beneficial. These
questions define the research gap addressed in this thesis.

\chapter{Method}

This research study was conducted at Neo4j, within the Graph Data Science team.
It follows an experimental systems-oriented methodology, where different GNN
training and inference pipelines are tested in an ablational manner. The
primary objective in the experiments is to find a GNN pipeline that keeps
functional equivalence to PyTorch Geometric while maximizing throughput.

\section{Experimental methodology}
The study is based around a graph-database-backed GNN pipeline that has been
proven to have functional equivalence to PyTorch Geometric. Components of the
GNN pipeline are modified and benchmarked to identify performance bottlenecks
and potential optimizations.

Each experiment isolates a single design choice, such as sampling strategy,
data retrieval method, or caching mechanism, allowing its impact on performance
to be evaluated independently. This enables a systematic exploration of the
design space of graph-database-backed GNN training.

\section{Baseline graph-database-backed GNN system}
As a reference, a PyTorch Geometric-based training pipeline is used. This
pipeline defines the expected behavior of sampling, batching, and model
execution, and serves as a ground truth for correctness. The Neo4j-backed
implementation replaces the data loading and sampling components of this
pipeline while keeping the model architecture and training loop identical.

\section{System configurations}
Experiments are conducted across multiple system configurations, where both
hardware setup and pipeline design may vary. These configurations include:

\begin{itemize}
    \item Different sampling implementations (Cypher queries vs. procedures)
    \item Variations in data retrieval strategies (batch size, serialization format)
    \item Use of caching mechanisms in feature retrieval
    \item Parallel vs. single-sampler execution
\end{itemize}

Each configuration represents a distinct point in the design space of
graph-database-backed GNN training.

\section{Datasets}
The datasets, consisting of homogeneous graphs, used for training and
evaluating GNNs for node classification were:

\begin{itemize}
    \item \textbf{ogbn-arxiv}, 170k nodes: a directed graph representing the citation network between computer science arXiv papers indexed by MAG.
    \item \textbf{ogbn-products}, 2.4M nodes: an undirected and unweighted graph representing an Amazon product co-purchasing network.
    \item \textbf{ogbn-papers100M}, 111M nodes: a directed citation graph of 111 million papers indexed by MAG.
\end{itemize}

\section{Equivalence constraint}
A key requirement throughout all experiments is that the Neo4j-backed pipeline
remains functionally equivalent to the PyTorch Geometric reference
implementation. Equivalence is evaluated using:

\begin{itemize}
    \item Consistency in sampled subgraph structure statistics, such as the number of sampled edges and nodes at each $K$-hop
    \item Similar training convergence behavior
    \item Comparable validation and test accuracy
    \item Distributional similarity of sampled neighborhoods at each $K$-hop
\end{itemize}

Only configurations that satisfy these criteria are considered valid, ensuring
that performance improvements do not come at the cost of altered model
behavior.

\section{Performance evaluation}
The primary performance metric is training throughput, measured as the number
of processed seed nodes per unit time. This is complemented by a detailed
breakdown of pipeline stages, including sampling time and feature retrieval
time. In addition, database-level metrics such as query latency, number of
database hits, and serialization overhead are collected to identify system
bottlenecks.

\section{Benchmarking procedure}
A dedicated measurement framework is used to collect timing and profiling data
during training. Each experiment is run for multiple epochs to ensure stable
measurements. Comparisons between configurations are made under identical
conditions, including a fixed model architecture, consistent hyperparameters,
and identical datasets. This ensures that observed performance differences are
attributable solely to system design choices.

\section{Graph neural network model}\label{sec:model}
All experiments use a two-layer Graph Convolutional Network (GCN). The model
applies two \texttt{GCNConv} message-passing layers with ReLU activations and a
final linear classifier. The architecture is intentionally kept simple to focus
evaluation on the data pipeline rather than the model itself. The model
parameters that vary are the dimension of the first \texttt{GCNConv} layer,
which matches the dimension of the node feature vectors, and the output
dimension, which matches the number of classes in the dataset.

\section{Sampling algorithms}\label{sec:sampling_algs}
The GraphSAGE sampler and the GraphSAINT sampler were implemented both in
Cypher and in Java using user-defined procedures in Neo4j. The sampling
implementations were then evaluated using the methodology described in this
chapter, maximizing throughput while keeping functional equivalence to the true
sampling algorithm. The code for all sampling implementations can be found in
the appendix.

\subsection{Neighbor sampling}
For each seed node in the mini-batch, the sampler issues a Cypher query that
traverses the graph up to $K$ hops, sampling at most
$[k_0, k_1, \ldots, k_{K-1}]$ neighbors at each hop. Neighbors exceeding the
limit are selected uniformly at random without replacement. This corresponds to
the GraphSAGE sampling strategy.

\begin{algorithm}[H]
\footnotesize
\caption{Neighbor Sampling}
\label{algo:neighbor_sampling}
\DontPrintSemicolon
\KwIn{Seed nodes $S$, neighbor limits $[k_0, k_1, \ldots, k_{K-1}]$}
\KwOut{Sampled node set $\mathcal{V}_{\mathrm{sub}}$}

$\mathcal{V}_{\mathrm{sub}} \leftarrow S$\;
$\mathrm{begin} \leftarrow 0$\;
$\mathrm{end} \leftarrow |S|$\;

\For{$\ell \leftarrow 0$ \KwTo $K-1$}{
    $k \leftarrow k_\ell$\;

    \For{$i \leftarrow \mathrm{begin}$ \KwTo $\mathrm{end}-1$}{
        $v \leftarrow \mathcal{V}_{\mathrm{sub}}[i]$\;
        Sample up to $k$ neighbors of $v$ uniformly at random, with or without replacement\;
        Add sampled neighbors to $\mathcal{V}_{\mathrm{sub}}$ if not already present\;
    }

    $\mathrm{begin} \leftarrow \mathrm{end}$\;
    $\mathrm{end} \leftarrow |\mathcal{V}_{\mathrm{sub}}|$\;
}
\Return{$\mathcal{V}_{\mathrm{sub}}$}\;
\end{algorithm}

\subsection{GraphSAINT random-walk sampling}
Instead of expanding neighborhoods around seed nodes, GraphSAINT constructs a
subgraph by first sampling a set of nodes via random walks and then inducing
the subgraph over those nodes.

\chapter{System Design}

The system consists of three main components: a Neo4j graph database acting as
the graph and feature store, a Python training process that implements the
PyTorch Geometric abstract interfaces, and a \texttt{NodeLoader} that
orchestrates mini-batch construction during training. The implementations of
\texttt{FeatureStore} and \texttt{GraphStore} act as the interface between the
training process and the Neo4j database. The implementation of
\texttt{BaseSampler} defines how the subgraph sampling is done.

\section{Baseline architecture}
The baseline GNN pipeline has three main components: (i) Neo4j to store the
graph, traverse the graph, and retrieve node feature vectors, node labels, and
data split information; and (ii) a Python process that implements the required
interfaces for abstract graph storage and also the model forward pass and
weight updates.

The Neo4j database is used to store both the graph structure and the features
of the nodes. The Python training process can therefore, instead of
materializing the full graph in memory, retrieve the data needed for only the
current mini-batch. For each batch of seed nodes, the Python process issues
queries to Neo4j through the implemented PyTorch Geometric interfaces.

Within the Python process, the \texttt{Neo4jGS} implementation is responsible
for graph traversal and neighborhood expansion, while \texttt{Neo4jAbstractFS}
handles feature retrieval. The \texttt{Neo4jSampler} coordinates these
components to produce sampled subgraphs that are functionally equivalent to
those expected by the standard PyTorch Geometric training pipeline. The
resulting architecture preserves the usual mini-batch GNN training loop, but
replaces in-memory graph access with on-demand retrieval from the database.

The baseline system is deployed on a single Google Cloud Compute Engine virtual
machine (\texttt{n2-highmem-32}) with 32 vCPUs and 256\,GB RAM. Both the Neo4j
database and the Python training process run on the same machine. This design
avoids inter-machine communication and makes it possible to isolate the
overhead introduced by database-backed sampling and feature retrieval.

\section{Improved architecture}
The improved architecture builds on the baseline design by optimizing the most
important parts of the data path: neighborhood sampling, serialization, feature
retrieval, and parallel batch preparation. The concrete improvements are
introduced and evaluated incrementally in the experiments rather than assumed
up front.

\section{Graph representation in Neo4j}\label{sec:graph_repr}
Nodes in Neo4j are stored with a label corresponding to the node type, for
example \texttt{Paper} for citation networks. Each node carries the following
properties:

\begin{itemize}
    \item A unique integer id.
    \item Pre-computed node feature vectors. The vectors can either be stored as \texttt{Float64[]} or \texttt{bytes[]}.
    \item A target class label for supervised node classification, represented as an integer.
    \item A split indicator, such as \texttt{train}, \texttt{val}, or \texttt{test}, used to partition the dataset.
\end{itemize}

Edges are stored as typed relationships connecting pairs of nodes. No edge
features are used in the experiments.

\section{PyTorch Geometric integration}\label{sec:pyg_impl}
Three PyTorch Geometric abstract classes described in Section~\ref{sec:pyg}
were implemented.

\noindent \textbf{\texttt{Neo4jAbstractFS}} is an abstract class extending
\texttt{FeatureStore}. The purpose of this class is to provide storage of the
feature vectors, for example in a graph database, and to implement an
efficient interface for fetching feature vectors. The reason it is kept
abstract is that while there are some methods any Neo4j-interacting PyTorch
Geometric \texttt{FeatureStore} needs, such as executing queries, there are
also some features that may vary depending on the use case. One example of
this is caching. Avoiding Neo4j fetches by storing commonly used features in
RAM is a good optimization, but how this caching is done in the feature store
should be up to the user. There are also two extensions of
\texttt{Neo4jAbstractFS}: \texttt{Neo4jCachedFS} and \texttt{Neo4jNoCacheFS}.

\noindent \textbf{\texttt{Neo4jAbstractGS}} is an abstract class extending
\texttt{GraphStore}, with the main goal of implementing efficient traversal
over the graph to enable fast subgraph sampling. It has two implementations,
\texttt{Neo4jMultiGS} and \texttt{Neo4jSingleGS}, where the difference is how
many samplers they are optimized to handle.

\noindent \textbf{\texttt{Neo4jSampler}} implements the sampling algorithm,
and is tightly coupled with the graph store where it executes the sampling
algorithm.

\section{Profiling and logging}
\subsection{Training logging}
Measurements are continuously taken during training and stored in a structured
format for later analysis.

\subsection{Query execution profiling}\label{sec:query_profiling}
Cypher queries are instrumented with the \texttt{PROFILE} keyword to collect
operator-level statistics: rows produced, database hits, estimated memory
usage, and query execution time. These are aggregated across batches and stored
as averaged CSV output.

\subsection{Python-level profiling}\label{sec:python_profiling}
Call-stack profiles are collected using Python's \texttt{cProfile} module and
stored as \texttt{train\_profile.txt}. These are used to identify hot spots in
the Python data pipeline that are not visible in the query-level metrics, such
as tensor allocation overhead or unnecessary data copies.

\section{Implementation decisions}
Important design tradeoffs in the system concern where to spend memory, where
to accept approximation, and how tightly to couple the PyTorch Geometric
interfaces to Neo4j-specific capabilities. Caching improves throughput but uses
more RAM. Procedures can outperform declarative queries for specialized
sampling, but increase implementation complexity and reduce portability.
Similarly, multi-sampler execution can improve concurrency, but it may also
increase contention in the database and complicate equivalence guarantees.

\chapter{Results}

This chapter presents the experimental results from graph-database-backed GNN
training and inference. The evaluation focuses on which phases dominate runtime
during training or inference, as well as total throughput and functional
correctness in the sampling.

\section{RQ1: Bottlenecks}
\subsection{Training phase vs. sampling phase}
\subsection{Sampling subphases}
\subsection{Feature vector datatype}
\subsection{Effect of caching}
\subsection{Effect of multisampler parallelism}
\subsection{Summary of RQ1}

\section{RQ2: Sampling strategies}
\subsection{Neighbor sampling in Neo4j}
\subsection{GraphSAINT in Neo4j}
\subsection{Summary of RQ2}

\section{RQ3: Inference with graph stored in Neo4j}
\subsection{Inference setting definitions}
\subsection{Batch inference performance}
\subsection{Inference on newly inserted nodes}
\subsection{Summary of RQ3}

\section{RQ4: Effect of graph size}
\subsection{Small graphs}
\subsection{Medium graphs}
\subsection{Large graphs}
\subsection{Summary of RQ4}

\chapter{Discussion}

This chapter discusses the meaning of the results in relation to the research
questions and the broader goal of assessing graph databases as a backend for
large-scale GNN training and inference.

\section{Optimal system implementation}
The discussion centers on which combination of sampling strategy, feature
retrieval method, caching, and execution model provides the best tradeoff
between throughput, simplicity, and correctness.

\section{Implications}
The implications are both practical and methodological. Practically, the work
helps clarify when database-backed GNN pipelines are a viable alternative to
in-memory or distributed systems. Methodologically, it contributes a way of
evaluating such systems under a strict equivalence constraint to a PyTorch
Geometric reference implementation.

\section{Threats to validity}
Threats to validity include the use of a single model family, a limited set of
datasets, hardware-specific effects, and implementation choices tied to Neo4j
and PyTorch Geometric. These constraints limit the generality of the results,
but they are acceptable given that the thesis focuses on a concrete systems
problem rather than universal model performance.

\chapter{Conclusion}

This thesis studies whether a graph database can act as an efficient and
scalable backend for GNN training and inference. By implementing a Neo4j-backed
pipeline integrated with PyTorch Geometric, the work investigates bottlenecks,
sampling efficiency, inference strategies, and how the usefulness of the
approach changes with graph size. The thesis contributes both a concrete system
design and an evaluation methodology for comparing graph-database-backed GNN
pipelines against a functionally equivalent reference implementation.

\section{Future work}
Future work includes broadening the evaluation to more GNN architectures,
testing additional graph databases, extending the system to more distributed
deployments, and studying how database-backed inference behaves under stricter
latency requirements and dynamic graph updates.

% Uncomment if needed.
% \appendix
% \chapter{Appendix Title}
% Supplementary material goes here.

% Uncomment if using BibTeX.
% \bibliographystyle{uclathes}
% \bibliography{references}

\end{document}
